\input{preamble.tex}

\usepackage{komorebi-monograph}
\renewcommand{\KomorebiKind}{Lean and evidence guide}
\renewcommand{\KomorebiShortTitle}{Paper III / Lean guide}
\renewcommand{\KomorebiCoverBlurb}{Read the exact source declarations, follow the map directions, and follow the current kernel verification and the preserved finite theorem map.}
\KomorebiSetup
\setmonofont{DejaVu Sans Mono}[Scale=.80]
\begin{document}
\KomorebiCover
\KomorebiOpeningNote{A practical guide to the current release and the complete 41-row finite algebraic contract. Read source-exact proof excerpts, follow the direction of each map, distinguish exact coordinate transport from independent connection data, and locate every formal target.
\medskip

\textbf{Evidence at this edition.} The release includes a fresh pinned compilation, a compiled ownership and axiom audit, upstream-only kernel replay, and adverse controls. The opening release chapter gives their exact scope. The preserved finite map still has 143 targets and its historical namespace audit has 468 entries. Those older counts describe the finite baseline, not the larger current closure.}


\tableofcontents
\clearpage
\pagenumbering{arabic}
\input{chapters/00_current_release.tex}
\clearpage
\input{chapters/01_orientation.tex}
\input{chapters/02_edges.tex}
\input{chapters/03_heterogeneous.tex}
\input{chapters/04_connections.tex}
\input{chapters/05_gauges.tex}
\input{chapters/06_nonlinear.tex}
\input{chapters/07_replay.tex}
\clearpage
\section{The complete 41-row formal contract}\label{sec:map}
This index preserves every row, target, status, and boundary of the authoritative v1.0.1 map. Fully qualified declaration names are printed in full. File and line locations are a current lexical navigation aid, not a second compiled type audit. Each target appears in exactly one row. The 143 names comprise 88 theorem declarations and 55 constructions.

The finite limits are part of the statements: edge sums use finite packet and edge types, arbitrary-length heterogeneous products keep one common finite packet type, and packet reconstruction uses a finite complete idempotent family. Rows 38--41 have the separate grammar and candidate-specific scopes explained in Section~6.
\input{generated/theorem_map.tex}
\clearpage
\appendix
\section{Monograph navigation}\label{app:crosswalk}
The following locations point into the monograph companion. Printed pages differ from a PDF viewer's physical page index because the cover is unnumbered and the opening pages use Roman numerals. A locator connects related presentations; it does not certify every neighboring ordinary statement in Lean.
\begin{longtable}{p{.54\textwidth}p{.36\textwidth}}
\toprule Canonical label & Final printed locator \\ \midrule\endhead
\code{sec:purpose} & Section 1, p.~1 \\ \addlinespace
\code{sec:edge-operators} & Section 2, p.~2 \\ \addlinespace
\code{def:edge-system} & Definition 2.1, p.~2 \\ \addlinespace
\code{lem:edge-support} & Lemma 2.2, p.~2 \\ \addlinespace
\code{thm:edge-annihilation} & Theorem 2.3, p.~2 \\ \addlinespace
\code{thm:edge-composable} & Theorem 2.4, p.~3 \\ \addlinespace
\code{sec:words-walks} & Section 3, p.~3 \\ \addlinespace
\code{def:word-op} & Definition 3.1, p.~3 \\ \addlinespace
\code{thm:first-break} & Theorem 3.2, p.~3 \\ \addlinespace
\code{thm:nonwalk-vanish} & Theorem 3.3, p.~4 \\ \addlinespace
\code{def:recursive-word-expansion} & Definition 3.4, p.~4 \\ \addlinespace
\code{thm:power-recursive-expansion} & Theorem 3.5, p.~4 \\ \addlinespace
\code{lem:word-endpoint} & Lemma 3.7, p.~5 \\ \addlinespace
\code{thm:word-to-walk} & Theorem 3.8, p.~6 \\ \addlinespace
\code{cor:word-to-walk-entry} & Corollary 3.9, p.~6 \\ \addlinespace
\code{cor:fixed-reachability} & Corollary 3.10, p.~6 \\ \addlinespace
\code{cor:orbit-reachability} & Corollary 3.11, p.~6 \\ \addlinespace
\code{sec:canonical-transport} & Section 4, p.~7 \\ \addlinespace
\code{thm:exact-frame-transport} & Theorem 4.1, p.~7 \\ \addlinespace
\code{thm:exact-path-transport} & Theorem 4.2, p.~7 \\ \addlinespace
\code{def:canonical-J} & Definition 4.3, p.~8 \\ \addlinespace
\code{thm:J-cocycle} & Theorem 4.4, p.~8 \\ \addlinespace
\code{thm:canonical-naturality} & Theorem 4.5, p.~9 \\ \addlinespace
\code{cor:J-flat} & Corollary 4.6, p.~9 \\ \addlinespace
\code{sec:connection} & Section 5, p.~9 \\ \addlinespace
\code{def:minimal-connection} & Definition 5.1, p.~9 \\ \addlinespace
\code{thm:minimal-flatness} & Theorem 5.2, p.~10 \\ \addlinespace
\code{prop:twisted-loop} & Proposition 5.3, p.~10 \\ \addlinespace
\code{def:connection} & Definition 5.5, p.~11 \\ \addlinespace
\code{def:calibration-flatness} & Definition 5.6, p.~12 \\ \addlinespace
\code{def:cocycle-defect} & Definition 5.8, p.~12 \\ \addlinespace
\code{def:triangle-holonomy} & Definition 5.9, p.~12 \\ \addlinespace
\code{thm:defect-holonomy} & Theorem 5.10, p.~13 \\ \addlinespace
\code{cor:canonical-zero-defect} & Corollary 5.11, p.~13 \\ \addlinespace
\code{def:path-transport} & Definition 5.12, p.~13 \\ \addlinespace
\code{prop:flat-endpoints} & Proposition 5.13, p.~14 \\ \addlinespace
\code{sec:gauge} & Section 6, p.~14 \\ \addlinespace
\code{def:gauge-change} & Definition 6.1, p.~14 \\ \addlinespace
\code{thm:gauge-covariance} & Theorem 6.2, p.~14 \\ \addlinespace
\code{thm:flat-gauge-canonical} & Theorem 6.3, p.~14 \\ \addlinespace
\code{sec:obstruction-packets} & Section 7, p.~15 \\ \addlinespace
\code{def:generic-obstruction} & Definition 7.1, p.~15 \\ \addlinespace
\code{thm:generic-obstruction} & Theorem 7.2, p.~15 \\ \addlinespace
\code{def:obstruction-packet} & Definition 7.4, p.~16 \\ \addlinespace
\code{thm:obstruction-detection} & Theorem 7.5, p.~17 \\ \addlinespace
\code{cor:packet-witness} & Corollary 7.6, p.~18 \\ \addlinespace
\code{prop:packet-gauge} & Proposition 7.7, p.~18 \\ \addlinespace
\code{sec:holonomy-walks} & Section 8, p.~19 \\ \addlinespace
\code{prop:heterogeneous-path} & Proposition 8.1, p.~19 \\ \addlinespace
\code{sec:example} & Section 9, p.~20 \\ \addlinespace
\code{sec:arbitrary-maps} & Section 10, p.~20 \\ \addlinespace
\code{prop:set-naturality} & Proposition 10.1, p.~20 \\ \addlinespace
\code{sec:operation-trees} & Section 11, p.~22 \\ \addlinespace
\code{def:compiled-operation-tree} & Definition 11.1, p.~22 \\ \addlinespace
\code{thm:tree-locality} & Theorem 11.2, p.~22 \\ \addlinespace
\code{def:signature} & Definition 11.4, p.~22 \\ \addlinespace
\code{def:operation-tree} & Definition 11.5, p.~23 \\ \addlinespace
\code{sec:ledger} & Section 12, p.~24 \\ \addlinespace
\bottomrule\end{longtable}

\section{Inherited declaration names}\label{app:inherited}
The following twelve names are inherited citations to the separately verified Papers I--II v1.2.1 bundle. The standalone Paper III build does not recompile them. Their fourteen manuscript uses are checked against the frozen manifest.
\input{generated/inherited.tex}
\clearpage
\section{Imports and exact dependency revisions}
The import table lists direct local file dependencies, not a theorem-level dependency graph. The adjacent JSON includes direct external imports as well. This appendix describes the preserved finite baseline. Its aggregate imports seven mathematical modules. The current release adds the continuation and packet calculus modules listed in the opening chapter.
\input{generated/imports.tex}
The manifest's exact pinned revisions follow. They accompany the toolchain string \code{leanprover/lean4:v4.34.0-rc2}.
\input{generated/pins.tex}
\clearpage
\section{Source and reproduction inventory}
The preserved finite baseline component has nine Lean files, totaling 1,948 source lines, and sixteen files including the map, manifest, gate, pins, and documentation. The guide stores 44 exact excerpts. For each, \code{generated/excerpts.json} records the source file, inclusive line range, full-source hash, exact-excerpt hash, and PDF-display-file hash.

\code{generated/theorem_map.json} contains all 41 rows and complete source locations. \code{generated/mapped_source_index.json} provides an alphabetical index. \code{generated/inherited_axiom_records.json} records the 143 saved transcript reports. \code{generated/static_audit.json} separates current checks from inherited counts and records the checksums used. The original report and transcript remain in the preserved finite release. Current verification sources and receipts are under \code{formal/} in the release root.

\code{build.py} is a compile-only XeLaTeX builder. \code{build_guide_data.py} regenerates static data from an explicitly supplied preserved release. \code{verify_pdf.py} renders the PDF, prepares contact sheets, and records page-level layout checks. \code{BUILD_NOTES.txt} records authoring failures and their repairs, rather than silently treating failed attempts as successful verification.

The guide uses P052 for body text, Nimbus Sans for headings, DejaVu Sans Mono for code, and Latin Modern Math for mathematical symbols. Code displays use \code{fvextra}, preserving token order. Display fallbacks are documented separately from the exact UTF-8 source excerpts.
\end{document}
